% Auto-generated by thesis/make_instruments.py --- do not edit by hand.
% Source: Novel Data/School Visit Proforma - *.docx

Two proformas were used on each visit. The lesson observation form was
completed once per lesson observed, and the outside-of-lessons form once per
school, covering transitions, break and lunch. Every item is rated 1--5 against
the anchor descriptions below, with anchors specified at 1, 3 and 5 and even
values used for intermediate judgements. Visits were made by teams of two or
three researchers, who completed the forms independently; item scores are
averaged across researchers and across lessons to give one value per school,
and \cref{sec:p1_irr_obs} reports agreement between observers.

\subsection*{Lesson observation}

\subsubsection*{Section 1 --- Strictness and Warmth}

\paragraph{Level of misbehaviour that results in a first sanction/level of misbehaviour which is tolerated}
Rated 1 (minor) to 5 (serious).

\begin{description}
\item[1] Isolated instances of misbehaviour such as whispering or fidgeting.
\item[3] Repeated disruptions, e.g., repeated talking or distraction of peers.
\item[5] Aggressive or defiant behaviour
\end{description}

\paragraph{Response of pupils to sanctions}
Rated 1 (very poor) to 5 (very good). A score of 0 is recorded separately (see below).

\begin{description}
\item[0] If no sanctions are observed
\item[1] Students ignore sanctions or react defiantly, showing no change in behaviour.
\item[3] Noticeable change in behaviour in response to sanctions, but not consistently maintained.
\item[5] Immediate and sustained positive response to sanctions e.g. apologises for actions.
\end{description}

\paragraph{Frequency of low-level disruption}
Rated 1 (very low) to 5 (very high).

\begin{description}
\item[1] Classroom is calm with disruptions rare or not occurring at all.
\item[3] Regular disruptions, sometimes affecting the flow of the lesson but generally manageable.
\item[5] Constant disruptions, creating a challenging environment for teaching and learning.
\end{description}

\paragraph{Concentration level of students}
Rated 1 (very low) to 5 (very high).

\begin{description}
\item[1] Students are consistently off-task, showing little interest or engagement with the lesson.
\item[3] Moderate levels of concentration, with students generally engaged but susceptible to distractions.
\item[5] Exceptional concentration, with students deeply engaged and consistently focused throughout the lesson.
\end{description}

\paragraph{Frequency of teacher-initiated participation}
Rated 1 (very low) to 5 (very high).

\begin{description}
\item[1] Teacher rarely encourages student participation, leading the class in a predominantly lecture-style format.
\item[3] Regular attempts by the teacher to involve students in the lesson, but some parts where teacher talks continuously.
\item[5] Teacher consistently involves students in the learning process.
\end{description}

\paragraph{Frequency of student-initiated participation}
Rated 1 (very low) to 5 (very high).

\begin{description}
\item[1] Students rarely or never take the initiative to participate, showing passivity in the learning process.
\item[3] Moderate frequency of student-initiated participation. Sometimes active but sometimes passive.
\item[5] Students consistently participate in the lesson demonstrating high levels of engagement in the learning process.
\end{description}

\paragraph{Motivation level of students}
Rated 1 (very low) to 5 (very high).

\begin{description}
\item[1] Students exhibit a lack of interest and motivation, showing little effort in classroom activities.
\item[3] A moderate level of motivation, with students showing interest at times but being disengaged at times.
\item[5] Exceptional motivation and enthusiasm, with students consistently demonstrating a strong desire to learn and participate.
\end{description}

\paragraph{Extent to which students address the teacher respectfully}
Rated 1 (not at all) to 5 (always).

\begin{description}
\item[1] Disrespectful interactions, use of inappropriate language or tone.
\item[3] Respectful interactions most of the time
\item[5] Exceptionally respectful and considerate interactions at all times.
\end{description}

\paragraph{Frequency with which teacher uses pupils' names}
Rated 1 (very low) to 5 (very high).

\begin{description}
\item[1] Rarely addresses students by name; impersonal interactions.
\item[3] Regularly uses student names, but some students are less frequently addressed.
\item[5] Consistently uses every student's name, fostering a sense of individual attention.
\end{description}

\paragraph{Frequency of praise}
Rated 1 (very low) to 5 (very high).

\begin{description}
\item[1] Teacher provides little to no praise, no positive reinforcement.
\item[3] Regular praise, but it is not always meaningful or specific.
\item[5] Regular and specific praise which highlights what students have done well.
\end{description}

\paragraph{Quality of interactions between teachers and pupils}
Rated 1 (impersonal) to 5 (warm).

\begin{description}
\item[1] Interactions are formal, distant, and lack warmth.
\item[3] Interactions are warm, but still quite formal in nature.
\item[5] Exceptionally warm and supportive interactions throughout the lesson.
\end{description}

\subsubsection*{Section 2 --- Teaching practice}

\paragraph{Frequency of questioning/checking pupil progress}
Rated 1 (very low) to 5 (very high).

\begin{description}
\item[1] Teacher rarely asks questions or checks on pupils' understanding.
\item[3] Teacher regularly asks standard questions, adequately assessing pupils' basic understanding but may not probe deeper comprehension.
\item[5] Teacher consistently integrates questioning seamlessly into lessons, skilfully using them to assess understanding and encourage higher-order thinking.
\end{description}

\paragraph{Quality of verbal feedback}
Rated 1 (very poor) to 5 (very good).

\begin{description}
\item[1] Feedback is rare, nonspecific, or irrelevant, offering no guidance or improvement for pupil learning.
\item[3] Feedback is regularly given, is clear, and somewhat constructive.
\item[5] Feedback is exceptionally well-tailored to each pupil and highly constructive.
\end{description}

\paragraph{Effectiveness of pupil discussion}
Rated 1 (not effective) to 5 (very effective). A score of 0 is recorded separately (see below).

\begin{description}
\item[0] No teacher-initiated pupil discussion is observed.
\item[1] Pupil discussions are rarely initiated or if they are, then they are poorly managed, resulting in minimal pupil engagement.
\item[3] Regularly encourages discussions which have a moderate impact on learning but may not fully engage all pupils.
\item[5] Regularly encourages discussions and structures them effectively to ensure that pupils engage with them and benefit from them.
\end{description}

\paragraph{Quality of differentiation (adaptation to students' needs)}
Rated 1 (very poor) to 5 (very good).

\begin{description}
\item[1] Teaching lacks adaptation to different learning needs, often leading to disengagement or unmet learning objectives for various students.
\item[3] Regularly employs differentiation strategies that address the needs of most students but not all.
\item[5] Exceptional in adapting teaching to meet all students' needs. Challenges high ability pupils and/or supports low ability pupils appropriately.
\end{description}

\paragraph{Clarity of teacher explanation}
Rated 1 (very poor) to 5 (very good).

\begin{description}
\item[1] Explanations are often unclear, leading to widespread confusion or misconceptions among students.
\item[3] Usually provides explanations that are clear and aid in student understanding, though some concepts may not be understood by all students.
\item[5] Provides exceptionally clear and comprehensive explanations, making complex concepts accessible to all or almost all students.
\end{description}

\paragraph{Extent to which learning outcomes are achieved}
Rated 1 (not at all) to 5 (completely).

\begin{description}
\item[1] Learning outcomes are rarely met, with most students failing to achieve the set objectives. Unclear what learning objectives are.
\item[3] A majority of learning outcomes are met, though some objectives may not be fully achieved by all students.
\item[5] All learning outcomes are fully met, with evidence of deep understanding and application of knowledge by all or almost all students.
\end{description}

\paragraph{Effectiveness of teaching methods}
Rated 1 (not effective at all) to 5 (very effective).

\begin{description}
\item[1] Teaching method(s) are not appropriate to the content and tasks.
\item[3] Teaching method(s) are mostly appropriate to the content and tasks.
\item[5] Chosen teaching method(s) are very well matched to the content and every task is clearly well thought out.
\end{description}

\paragraph{Clarity of lesson structure}
Rated 1 (unclear) to 5 (very clear).

\begin{description}
\item[1] Lessons lack a clear structure, resulting in disorganized content delivery. Structure is not communicated to students.
\item[3] Generally clear lesson structure with a logical flow, though minor issues in transition or clarity may occasionally arise. Structure is communicated quite well to students.
\item[5] Lessons are exceptionally well-structured, with seamless transitions, and a well-thought-out progression. Structure is regularly communicated to students.
\end{description}

\paragraph{Effectiveness of resource use}
Rated 1 (low) to 5 (high).

\begin{description}
\item[1] The resources do not allow most students to achieve the learning objectives and/or engage with the topic.
\item[3] The resources allow most students to achieve the learning objectives and/or engage with the topic.
\item[5] The resources allow all or almost all of the students to achieve the learning objectives and/or engage with the topic.
\end{description}


\subsection*{Outside-of-lessons observation}

\subsubsection*{Section 1 --- Comparison across lessons}

\paragraph{Consistency of use of sanction system}
Rated 1 (not consistent at all) to 5 (completely consistent).

\begin{description}
\item[1] Teachers are using different sanction systems in every classroom.
\item[3] Teachers are generally taking the same approach to misbehaviour, but there is no shared language.
\item[5] All teachers use the sanction system in exactly the same way and there is a clear shared language.
\end{description}

\paragraph{Consistency of use of reward system}
Rated 1 (not consistent at all) to 5 (completely consistent).

\begin{description}
\item[1] Teachers are using different reward systems in every classroom.
\item[3] Teachers are generally taking the same approach to positive behaviours, but there is no shared language.
\item[5] All teachers use the reward system in exactly the same way and there is a clear shared language.
\end{description}

\paragraph{Consistency of verbal feedback}
Rated 1 (not consistent at all) to 5 (completely consistent).

Should be calculated as the average of the scores from individual lessons.

\paragraph{Quality of differentiation (adaptation to students' needs) across lessons}
Rated 1 (very poor) to 5 (very good).

Should be calculated as the average of the scores from individual lessons.

\subsubsection*{Section 2 --- Transitions}

\paragraph{Calmness of corridors}
Rated 1 (chaotic) to 5 (calm).

\begin{description}
\item[1] The corridors feel unsafe. There is a lot of boisterous behaviour and students do not seem concerned with getting to their next lesson.
\item[3] Most students move between lessons in a calm fashion. There is an occasional loud voice and a small amount of boisterous behaviour. Most students move purposefully to their next lesson.
\item[5] All students move between lessons in a calm fashion with no loud voices or boisterous behaviour at all. All students move with purpose to their next lesson.
\end{description}

\paragraph{Arrival time to next lesson}
Rated 1 (lots of lateness) to 5 (all on time).

\begin{description}
\item[1] The majority of students arrive late to lessons.
\item[3] Most students arrive on time with a few late arrivals.
\item[5] All or almost all students arrive on time.
\end{description}

\paragraph{Quality of relationships between students during transition}
Rated 1 (very poor) to 5 (very good).

\begin{description}
\item[1] Lots of students speak to each other disrespectfully. There is clear evidence of bullying and/or intimidation.
\item[3] Most students are polite in their interactions with each other. There is evidence that a small amount of bullying/intimidation is taking place.
\item[5] All or almost all students are polite in their interactions with each other. There is no evidence of bullying or intimidation of any kind.
\end{description}

\paragraph{Presentation of wall/corridor displays}
Rated 1 (poor condition) to 5 (excellent condition).

\begin{description}
\item[1] There are no displays, or if there are, then they are old and tattered.
\item[3] Most displays are in good condition, but some are a bit tattered and look old.
\item[5] All displays are in pristine condition with no damage. They all look as if they are new.
\end{description}

\paragraph{Extent to which wall/corridor displays which are not student work align with school values}
Rated 1 (random) to 5 (complete alignment).

\begin{description}
\item[1] There are no displays which convey messages about the school values.
\item[3] There are some displays which convey messages about the school values.
\item[5] All or almost all displays convey messages about the school values.
\end{description}

\subsubsection*{Section 3 --- Break and lunch}

\paragraph{Quality of non-behaviour related interactions between staff and students}
Rated 1 (very poor) to 5 (very good).

\begin{description}
\item[1] In general, staff do not engage with students.
\item[3] Most staff actively engage with students and interactions are mostly warm and personal.
\item[5] All or almost all staff actively engage with students and interactions are warm and personal.
\end{description}

\paragraph{Quality of relationships between students}
Rated 1 (very poor) to 5 (very good).

\begin{description}
\item[1] Lots of students speak to each other disrespectfully. There is clear evidence of bullying and/or intimidation.
\item[3] Most students are polite in their interactions with each other. There is evidence that a small amount of bullying/intimidation is taking place.
\item[5] All or almost all students are polite in their interactions with each other. There is no evidence of bullying or intimidation of any kind.
\end{description}

\paragraph{Organisation of canteen/dining hall}
Rated 1 (chaotic) to 5 (calm).

\begin{description}
\item[1] Students get their food in a haphazard fashion. They leave food behind on tables and on the floor. They do not take care of their environment.
\item[3] There is some organisation to the way that students get their food. There is a little bit of mess left behind.
\item[5] Students get their food in an organised fashion. They leave the table and floor where they eat clean and tidy.
\end{description}

\paragraph{Organisation of recreational space}
Rated 1 (chaotic) to 5 (calm).

\begin{description}
\item[1] Activities are interfering with each other and there is no supervision from staff. There is a general sense of disorganisation. There is lots of rule breaking.
\item[3] There are clear areas for different activities, and some are supervised by staff. There is some rule breaking.
\item[5] There are clear areas for different activities, and these are all supervised by members of staff. All or almost all students abide by the rules.
\end{description}
